\documentclass[10pt,a4paper]{amsart}
\usepackage[margin=19mm]{geometry}
\usepackage{amsmath,amssymb,mathtools}
\usepackage[scr=rsfs,cal=euler]{mathalfa}
\usepackage{xfrac}
\usepackage[unicode,hidelinks,bookmarks=false]{hyperref}
\newcommand{\ie}{i.e.\ }
\newcommand{\cf}{cf.\ }
\newcommand{\resp}{respectively\ }
\newcommand{\Subsection}{\subsection}
\providecommand{\nobreakdash}{\nobreak\hbox{-}}
\setlength{\emergencystretch}{2em}
\multlinegap0pt
\usepackage{bm}
\usepackage{bbm}
\usepackage{dsfont}
\usepackage{xspace}
\usepackage{cancel}
\usepackage{mathtools}
\usepackage{amsthm}
\makeatletter
\@ifundefined{theorem}{\newtheorem{theorem}{Theorem}}{}
\@ifundefined{proposition}{\newtheorem{proposition}[theorem]{Proposition}}{}
\makeatother
\title[Reversibility and symmetry of affine toral automorphisms]{Reversibility and symmetry of affine toral automorphisms}
\author{Kuntal Banerjee, Anubrato Bhattacharyya, Krishnendu Gongopadhyay,, Subhamoy Mondal}
\date{Source: arXiv:2601.15827v2 (CC BY 4.0)}
\begin{document}
\maketitle

\noindent\emph{Source and licence.} Reversibility and symmetry of affine toral automorphisms. Version arXiv:2601.15827v2 (CC BY 4.0), distributed under the Creative Commons Attribution 4.0 International licence, as recorded in the arXiv licence metadata for this exact version and shown on the abstract page https://arxiv.org/abs/2601.15827v2. Full rights notice: licenses/arxiv-2601.15827-CC-BY-4.0.txt.\par\smallskip

\noindent\textit{Editorial scope.} Counted unit: the Theorem whose source label is \texttt{top entrpy} in the article (source0.tex lines 439--445, source label \texttt{top entrpy}), delivered together with the article's own complete original proof of it (source0.tex lines 446--505, one \texttt{proof} environment of 60 lines). No mathematics is written, repaired or re-derived: statement and proof are the article's own text line for line. Required context retained uncounted: conjugacy (lines 400--403). Every definition, display and label of those lines is retained unchanged. Omitted from this package and not redistributed: 1--399, 404--438, 506--1258. Nothing inside a retained excerpt is omitted or altered: every retained body is delivered exactly as the article's lines are written, so no omission receipt of any kind accompanies this package. Citations: the retained excerpts carry no citation command at all, so no bibliography entry of the article is transcribed and no reference is redistributed. Reference adaptation: none was needed and none was made; every reference inside the delivered unit resolves inside it, so no unresolved reference or citation is produced. Printed slips kept literal: no printed word, symbol, hypothesis or number of the counted statement or of its proof was altered. Layout and class: the article's own class is \texttt{amsart} with packages including geometry, booktabs, amssymb, xcolor, epstopdf, graphicx, tikz-cd, todonotes, ulem, subfig, natbib, tabulary; the wrapper is the lane's pinned \texttt{amsart} editorial wrapper at 10pt on A4, which restates the theorem-like environments the retained text needs and adds the article's own macro definitions that the retained text needs (none), with extra wrapper packages (bm, bbm, dsfont, xspace, cancel, mathtools). The delivered numbering is the unit's own. Assets: no retained span contains a figure environment, a graphics-inclusion command, a PDF-page inclusion, a table or a TikZ picture; no image file is copied into this package, the package declares no asset dependency and no image file is redistributed with it. Licence: version arXiv:2601.15827v2 of this article is distributed under the Creative Commons Attribution 4.0 International licence, as recorded in the arXiv licence metadata for this exact version and shown on the abstract page https://arxiv.org/abs/2601.15827v2. A full rights notice is delivered at licenses/arxiv-2601.15827-CC-BY-4.0.txt. Characters: every retained excerpt is pure ASCII, so the delivered engine, XeLaTeX with its default Latin Modern text font, reports no missing character.
\begin{proposition}\label{conjugacy}
Let $A \in GL(2,\mathbb{Z})$ with $1$ not an eigenvalue.  
Then the maps $f_{A,\overline{b}}: \bar{x} \mapsto A\bar{x} + \overline{b}~(\text{mod}~\mathbb{Z}^2)$ and $f_A:\bar{x} \mapsto A\bar{x}~(\text{mod}~\mathbb{Z}^2)$ are topologically conjugate on $\mathbb{T}^2$ via a translation.
\end{proposition}

\begin{theorem}\label{top entrpy}
   Let $A \in GL(2,\mathbb{Z})$ and $\bar{a} \in \mathbb{T}^2$. Then the toral automorphism $f_A$ and the affine toral automorphism $f_{A,\bar{a}}$ have the same topological entropy; that is,
\[
h_{\mathrm{top}}(f_A) = h_{\mathrm{top}}(f_{A,\bar{a}}).
\]
 
\end{theorem}

\begin{proof}
If $1$ is not an eigen value of $A,$ then the proof follows from Proposition \ref{conjugacy}.    

Now assume 1 is an eigenvalue of $A \in GL(2,\mathbb{Z})$. This means the other eigenvalue is either $1$ or $-1$. If both eigenvalues of $A$ are equal to $1$ then there exists an invertible matrix $P$ so that, $PAP^{-1}=I+N, ~N^2=O \in M(2, \mathbb Z)$ and hence $PA^nP^{-1}=I+nN \implies ||PA^nP^{-1}|| \leq 1+n||N||=1+n$, where $||.||$ denotes the operator norm of a matrix. This means $d(f_A^n(x),f_A^n(y))\leq nC \cdot d(x,y)$ ($C$ is independent of $n$), where $d$ is the induced metric on the the torus.

Recall the Bowen metric $d_n$ on a compact metric space $X$, which measures the distance between orbits of length $n$,  $$d_n(x, y) = \max_{0 \leq i < n} d(f_A^i(x), f_A^i(y)).$$ Using the Lipschitz condition on each iterate $f_A^i$, we have,  $d(f_A^i(x), f_A^i(y))  \leq iC  d(x, y)$ and hence for $i \leq n$, we get that $d_n(x, y) \leq nC  d(x, y)$.

Let $N_n(\mathbb{T}^2, \epsilon)$ be the minimum number of $\epsilon$-balls in the $d_n$ metric required to cover $\mathbb{T}^2$. From the inequality above, we see that a ball in the original metric $d$ with radius $ \frac{\epsilon}{Cn}$ is entirely contained within $\epsilon$-ball in the $d_n$ metric. Therefore, the number of balls needed to cover $\mathbb{T}^2$ in the $d_n$ metric is at most the number of balls needed to cover $\mathbb{T}^2$ in the original metric with radius $\frac{\epsilon}{Cn}$ i.e., $N_n(\mathbb{T}^2, \epsilon) \leq N\left(\mathbb{T}^2, \frac{\epsilon}{Cn}\right),$ where $N(\mathbb{T}^2, \delta)$ is the minimum number of balls of radius $\delta$ required to completely cover   $\mathbb{T}^2$ in the $d$ metric. Clearly, $N(\mathbb{T}^2, \delta)=O(1/\delta^2)$, which means, 
\[
\begin{array}{c}
N_n(\mathbb{T}^2, \epsilon)
\leq N\!\left(\mathbb{T}^2, \dfrac{\epsilon}{Cn}\right)
\leq \dfrac{C_0 n^2}{\epsilon^2}. \\[0.6em]
\end{array}
\]

This yields, 

\[
\begin{array}{c}
\dfrac{\log N_n(\mathbb{T}^2, \epsilon)}{n}
\leq
\dfrac{\log C_0 + 2\log n - 2\log \epsilon}{n}
\end{array}
\]

hence,  
\[
\begin{array}{rcl}
\displaystyle
\limsup_{n \to \infty} \frac{1}{n}\log N_n(\mathbb{T}^2,\epsilon)
& \le &
\displaystyle
\limsup_{n \to \infty} \frac{\log C_0 + 2\log n - 2\log \epsilon}{n} \\[0.8em]
& = &
\displaystyle
\lim_{n \to \infty} \frac{\log C_0 + 2\log n - 2\log \epsilon}{n} \\[0.8em]
& = & 0 .
\end{array}
\]

Therefore, the topological entropy of the toral automorphism $f_A$ is
\[
h_{\mathrm{top}}(f_A)
=
\lim\limits_{\epsilon \to 0}
\limsup_{n \to \infty}
\frac{1}{n}\log N_n(\mathbb{T}^2,\epsilon)
= 0 .
\]

If the other eigenvalue is $-1$ then $A$ is conjugate to 
$\begin{pmatrix}
1 & 0\\
0 & -1
\end{pmatrix}$ so that $||A^n||\leq C$. It is clear that the same argument as above applies to this case. Hence $h_{\mathrm{top}}(f_{A})=0$, if one of the eigenvalues is equal to $1$. 

Since $f_{A,\bar{a}}$ is just a translation of $f_A$ by $\bar{a}$, the Lipschitz bound on $d_n$ metric remains unchanged and the whole proof above applies. Hence, we can conclude that $h_{\mathrm{top}}(f_A)=h_{\mathrm{top}}(f_{A,\bar{a}})=0$, if one of the eigenvalues is equal to 1.

\end{proof}

\end{document}
